% ============================================================
\chapter{Diffusion and Transport Models}
\label{ch:diffusion}
% ============================================================

\section{Introduction}

Diffusion is everywhere: a dye disperses in water, heat spreads through a metal, a pollutant disseminates in the atmosphere. This fundamental phenomenon, first studied quantitatively by Adolf Fick in 1855 (by analogy with Fourier's law for heat), obeys the diffusion equation---the same heat equation, but interpreted in terms of concentration rather than temperature. Coupled with a velocity field, it becomes the advection-diffusion equation, which models the transport of substances in a moving medium. This chapter develops these models and their applications.

Diffusion and transport phenomena are ubiquitous: heat conduction,
pollutant dispersion, cell migration.  This chapter derives the heat
equation from Fick's law, studies the advection-diffusion equation,
and introduces numerical methods for parabolic PDEs.

% ------------------------------------------------------------
\section{Fick's Law and the Diffusion Equation}
\label{sec:fick}
% ------------------------------------------------------------

\begin{definition}[Fick's law]
The diffusive flux $J$ of a substance with concentration $c(x,t)$ is
proportional and opposite to the concentration gradient:
\[
  J = -D\frac{\partial c}{\partial x},
\]
where $D>0$ is the \textbf{diffusion coefficient} (m$^2$/s).
\end{definition}

\begin{theorem}[Diffusion (heat) equation]
Combining Fick's law with mass conservation
$\partial c/\partial t = -\partial J/\partial x$ yields:
\[
  \frac{\partial c}{\partial t} = D\frac{\partial^2 c}{\partial x^2}.
\]
This is the one-dimensional heat equation.
\end{theorem}

\begin{proof}
Local conservation:
\[
  \frac{\partial c}{\partial t} = -\frac{\partial J}{\partial x}
  = -\frac{\partial}{\partial x}\!\left(-D\frac{\partial c}{\partial x}\right)
  = D\frac{\partial^2 c}{\partial x^2}.
\]
\end{proof}

\begin{keyformulas}[title=\textbf{Fundamental solution (heat kernel)}]
For the problem on $\R$ with initial condition $c(x,0)=\delta(x)$:
\[
  c(x,t) = \frac{1}{\sqrt{4\pi Dt}}\,\exp\!\left(-\frac{x^2}{4Dt}\right),
  \qquad t>0.
\]
This is a Gaussian whose standard deviation grows as $\sigma(t)=\sqrt{2Dt}$.
\end{keyformulas}

\begin{intuition}[title=\textbf{Diffusion and random walks}]
Diffusion is the macroscopic limit of a random walk.  If a particle
takes steps of length $\ell$ at rate $1/\Delta t$, the diffusion
coefficient is $D=\ell^2/(2\Delta t)$.
\end{intuition}

% ------------------------------------------------------------
\section{Boundary Conditions}
% ------------------------------------------------------------

\begin{definition}[Classical boundary conditions]
On a domain $[0,L]$, the most common boundary conditions are:
\begin{itemize}
  \item \textbf{Dirichlet:} $c(0,t)=c_0$, $c(L,t)=c_L$ (prescribed
    concentration).
  \item \textbf{Neumann:} $\partial c/\partial x(0,t)=0$ (zero flux,
    insulating wall).
  \item \textbf{Robin:} $-D\partial c/\partial x + hc = hc_\infty$
    (convective transfer).
\end{itemize}
\end{definition}

% ------------------------------------------------------------
\section{Finite Difference Methods}
\label{sec:finite_diff}
% ------------------------------------------------------------

\begin{definition}[Discretisation]
Discretise $[0,L]$ into $N+1$ points $x_j=j\Delta x$ and time into
steps $\Delta t$.  The second derivative is approximated by:
\[
  \frac{\partial^2 c}{\partial x^2}\bigg|_{x_j}
  \approx \frac{c_{j-1}-2c_j+c_{j+1}}{(\Delta x)^2}.
\]
\end{definition}

\begin{theorem}[Explicit scheme (FTCS)]
The \textit{Forward Time, Central Space} scheme is:
\[
  c_j^{n+1} = c_j^n + \mu(c_{j-1}^n - 2c_j^n + c_{j+1}^n),
  \qquad \mu = \frac{D\Delta t}{(\Delta x)^2}.
\]
This scheme is stable if and only if $\mu\leq 1/2$.
\end{theorem}

\begin{warning}[title=\textbf{CFL condition}]
The stability condition $\mu = D\Delta t/(\Delta x)^2 \leq 1/2$
links the time and space steps.  For fine meshes, $\Delta t$ must
be very small, making the explicit scheme costly.  Implicit schemes
(Crank--Nicolson) remove this constraint.
\end{warning}

\begin{codeblock}[title=\textbf{1D Diffusion --- explicit scheme}]
\begin{minted}{python}
import numpy as np
import matplotlib.pyplot as plt

L, D, T_final = 1.0, 0.01, 2.0
Nx = 50; dx = L / Nx
dt = 0.4 * dx**2 / D  # CFL: mu = 0.4 < 0.5
Nt = int(T_final / dt)
mu = D * dt / dx**2

x = np.linspace(0, L, Nx + 1)
c = np.exp(-((x - 0.5)**2) / (2 * 0.05**2))

plt.figure(figsize=(10, 6))
plt.plot(x, c, label='$t=0$')

for n in range(Nt):
    c_new = c.copy()
    c_new[1:-1] = c[1:-1] + mu * (c[:-2] - 2*c[1:-1] + c[2:])
    c_new[0] = 0; c_new[-1] = 0
    c = c_new
    if (n+1) * dt in [0.05, 0.2, 0.5, 1.0, 2.0]:
        plt.plot(x, c, label=f'$t={((n+1)*dt):.2f}$')

plt.xlabel('$x$'); plt.ylabel('$c(x,t)$')
plt.legend(); plt.title('1D Diffusion (explicit scheme)')
plt.tight_layout(); plt.savefig('diffusion_explicit.pdf')
\end{minted}
\end{codeblock}

% ------------------------------------------------------------
\section{The Advection-Diffusion Equation}
\label{sec:advection_diffusion}
% ------------------------------------------------------------

\begin{definition}[Advection-diffusion equation]
When the substance is also transported by a velocity field $v$:
\[
  \frac{\partial c}{\partial t} + v\frac{\partial c}{\partial x}
  = D\frac{\partial^2 c}{\partial x^2}.
\]
The P\'eclet number $\mathrm{Pe}=vL/D$ measures the relative
importance of advection to diffusion.
\end{definition}

\begin{proposition}[Limiting regimes]
\begin{itemize}
  \item $\mathrm{Pe}\ll 1$: diffusion dominates (heat equation).
  \item $\mathrm{Pe}\gg 1$: advection dominates (transport equation).
  \item $\mathrm{Pe}\sim 1$: both effects are comparable.
\end{itemize}
\end{proposition}

\begin{codeblock}[title=\textbf{Advection-diffusion}]
\begin{minted}{python}
import numpy as np
import matplotlib.pyplot as plt

L, D, v = 1.0, 0.005, 0.5
Nx = 200; dx = L / Nx
dt = 0.001; Nt = 1000

x = np.linspace(0, L, Nx + 1)
c = np.exp(-((x - 0.2)**2) / (2 * 0.03**2))

plt.figure(figsize=(10, 5))
plt.plot(x, c, 'k--', label='$t=0$')

for n in range(Nt):
    c_new = c.copy()
    c_new[1:-1] = (c[1:-1]
                   - v * dt / dx * (c[1:-1] - c[:-2])
                   + D * dt / dx**2 * (c[:-2] - 2*c[1:-1] + c[2:]))
    c_new[0] = 0; c_new[-1] = 0
    c = c_new
    if (n+1) in [200, 500, 800, 1000]:
        plt.plot(x, c, label=f'$t={(n+1)*dt:.2f}$')

plt.xlabel('$x$'); plt.ylabel('$c(x,t)$')
plt.legend(); plt.title(f'Advection-diffusion (Pe = {v*L/D:.0f})')
plt.tight_layout(); plt.savefig('advection_diffusion.pdf')
\end{minted}
\end{codeblock}

% ------------------------------------------------------------
\section{Diffusion in Higher Dimensions}
% ------------------------------------------------------------

\begin{definition}[Diffusion equation in 2D/3D]
In $d$ spatial dimensions:
\[
  \frac{\partial c}{\partial t} = D\nabla^2 c
  = D\sum_{i=1}^d \frac{\partial^2 c}{\partial x_i^2}.
\]
The fundamental solution in $d$ dimensions is:
\[
  c(\mathbf{x},t) = \frac{1}{(4\pi Dt)^{d/2}}
  \exp\!\left(-\frac{\norm{\mathbf{x}}^2}{4Dt}\right).
\]
\end{definition}

% ------------------------------------------------------------
\section{Reaction-Diffusion Models}
\label{sec:reaction_diffusion}
% ------------------------------------------------------------

\begin{definition}[Reaction-diffusion equation]
Adding a reaction term $f(c)$:
\[
  \frac{\partial c}{\partial t} = D\frac{\partial^2 c}{\partial x^2}
  + f(c).
\]
\end{definition}

\begin{example}[Fisher--KPP equation]
With $f(c)=rc(1-c)$ (logistic growth), we obtain the Fisher equation:
\[
  \frac{\partial c}{\partial t} = D\frac{\partial^2 c}{\partial x^2}
  + rc(1-c).
\]
It admits \textbf{travelling wave fronts} with minimum speed
$v_{\min}=2\sqrt{Dr}$.
\end{example}

\begin{theorem}[Minimum front speed for Fisher--KPP]
The Fisher--KPP equation admits travelling wave solutions
$c(x,t)=\phi(x-vt)$ for all speeds $v\geq v_{\min}=2\sqrt{Dr}$.
For compactly supported initial data, the front propagates at speed
$v_{\min}$.
\end{theorem}

\begin{codeblock}[title=\textbf{Fisher wave front}]
\begin{minted}{python}
import numpy as np
import matplotlib.pyplot as plt

L, D, r = 20.0, 1.0, 1.0
Nx = 400; dx = L / Nx
dt = 0.4 * dx**2 / D
x = np.linspace(0, L, Nx + 1)

c = np.where(x < 2, 1.0, 0.0)

plt.figure(figsize=(10, 5))
times_to_plot = [0, 2, 4, 6, 8]
t_current = 0

for t_target in times_to_plot:
    while t_current < t_target:
        c_new = c.copy()
        c_new[1:-1] = (c[1:-1]
            + dt * D * (c[:-2] - 2*c[1:-1] + c[2:]) / dx**2
            + dt * r * c[1:-1] * (1 - c[1:-1]))
        c_new[0] = 1.0; c_new[-1] = 0.0
        c = c_new
        t_current += dt
    plt.plot(x, c, label=f'$t = {t_target}$')

v_min = 2 * np.sqrt(D * r)
plt.xlabel('$x$'); plt.ylabel('$c(x,t)$')
plt.title(f'Fisher Wave Front ($v_{{\\min}} = {v_min:.2f}$)')
plt.legend(); plt.tight_layout()
plt.savefig('fisher_front.pdf')
\end{minted}
\end{codeblock}

% ------------------------------------------------------------
\section{Turing Patterns}
% ------------------------------------------------------------

\begin{definition}[Turing instability]
A two-species reaction-diffusion system:
\begin{align}
  \frac{\partial u}{\partial t} &= D_u\nabla^2 u + f(u,v), \\
  \frac{\partial v}{\partial t} &= D_v\nabla^2 v + g(u,v),
\end{align}
may exhibit a \textbf{Turing instability}: a spatially homogeneous
equilibrium that is stable without diffusion can become unstable with
diffusion if $D_v \gg D_u$.  This gives rise to \textbf{spatial
patterns} (spots, stripes).
\end{definition}

\begin{remark}
Turing instability typically requires an activator (slow diffusion,
small $D_u$) and an inhibitor (fast diffusion, large $D_v$).
\end{remark}

% ------------------------------------------------------------
\section{Exercises}
% ------------------------------------------------------------

\begin{exercise}
Verify that the heat kernel
$c(x,t)=\frac{1}{\sqrt{4\pi Dt}}\exp(-x^2/(4Dt))$ satisfies the
diffusion equation.
\end{exercise}

\begin{exercise}
Solve the heat equation on $[0,L]$ with homogeneous Dirichlet
boundary conditions by separation of variables.  Show that the
general solution is:
\[
  c(x,t) = \sum_{n=1}^\infty B_n \sin\!\left(\frac{n\pi x}{L}\right)
  \exp\!\left(-\frac{n^2\pi^2 D}{L^2}t\right).
\]
\end{exercise}

\begin{exercise}
Implement the Crank--Nicolson scheme for 1D diffusion and compare
with the explicit scheme in terms of accuracy and stability.
\end{exercise}

\begin{exercise}
Simulate the advection-diffusion equation for different values of Pe
and observe the transition between the diffusive and advective regimes.
\end{exercise}

\begin{exercise}
For the Fisher--KPP equation, verify numerically that the front speed
converges to $v_{\min}=2\sqrt{Dr}$ for compactly supported initial
data.
\end{exercise}

\begin{exercise}
Implement a Schnakenberg-type reaction-diffusion system in 2D and
observe the formation of Turing patterns.
\end{exercise}
